\documentclass[12pt, a4paper]{article}

\usepackage{amsmath, amssymb, amsfonts}
\usepackage{float}
\usepackage{graphicx}
\usepackage{listings}
\usepackage{xcolor}
\usepackage[linesnumbered, ruled]{algorithm2e}
\usepackage{subcaption}
\usepackage[hidelinks]{hyperref}
% Each algorithm is written twice: the algorithm2e block (what the PDF typesets)
% and a plain-text Verbatim block right after it (what the web shows — pandoc
% renders Verbatim as a copyable <pre> code block). \excludecomment makes pdflatex
% SKIP the Verbatim blocks, so the PDF shows only the algorithm2e version; pandoc
% conversely drops each algorithm2e to an (empty, CSS-hidden) div and keeps the
% Verbatim. Because pdflatex never sees the Verbatim, its Unicode maths glyphs
% (⟨⟩ ← τ ∈ ∞ …) need no inputenc/literate setup.
\usepackage{comment}
\excludecomment{Verbatim}

\definecolor{codegreen}{rgb}{0,0.6,0}
\definecolor{codegray}{rgb}{0.5,0.5,0.5}
\definecolor{codepurple}{rgb}{0.58,0,0.82}
\hypersetup{linktoc=all}

\lstset{
    language=Python,
    basicstyle=\ttfamily\small,
    keywordstyle=\color{blue},
    commentstyle=\color{codegreen},
    stringstyle=\color{codepurple},
    breaklines=true,
    numbers=left,
    numberstyle=\tiny\color{codegray},
    frame=single,
    showstringspaces=false
}

\title{AI Masterclass Test 1}
\author{}
\date{}


\begin{document}

\section{Depth First Search}
\begin{algorithm}[H]
% \caption{DFS}
\SetKwFunction{Fdfs}{dfs}
\SetKwProg{Fn}{function}{}{end function}
\Fn{\Fdfs{$\mathcal{P} = \langle S, s_0, A, \tau, c, G \rangle$}}{
    $agenda \leftarrow newStack()$\;
    $agenda.push(s_0)$\;
    \While{not agenda.isEmpty()}{
        $s \leftarrow agenda.pop()$\;
        \tcp{expand s}
        \For{$a \in A$}{
            $s' \leftarrow \tau(s, a)$\;
            \If{$s' \in G$}{
                \tcp{goal found?}
                \Return "done"\;
            }
            \tcp{goal not found so add s' to stack for evaluation}
            $agenda.push(s')$\;
        }
    }
}
\end{algorithm}
\begin{Verbatim}
function dfs(P = ⟨S, s₀, A, τ, c, G⟩):
    agenda ← newStack()
    agenda.push(s₀)
    while not agenda.isEmpty():
        s ← agenda.pop()              # expand s
        for a ∈ A:
            s' ← τ(s, a)
            if s' ∈ G:                # goal found?
                return "done"
            agenda.push(s')           # goal not found, add s' to stack
\end{Verbatim}

\begin{itemize}
    \item Utilise stack for LIFO
    \item Expands one branch at a time till the deepest node
\end{itemize}

\noindent Evaluation
\begin{itemize}
    \item Completeness: Not guaranteed to find a solution if one exists
    \item Time complexity: If it finds a solution, it takes much less time than BFS
    \item Space complexity: Current branch - $d$ nodes at depth $d$.
    \item Optimality: Solution found not guaranteed to be optimal.
\end{itemize}

\end{document}
